\chapter*{Resumen}
\label{resumen}
En este trabajo se aborda el problema de la congestión vehicular en áreas urbanas mediante un sistema de control semafórico basado en el \gls{acr-marl} cooperativo. Se analiza la evolución del parque automotor en Argentina, se describen las limitaciones del control semafórico tradicional y se formula el problema como una tarea de decisión secuencial en la que las intersecciones actúan como agentes. El objetivo es desarrollar y evaluar políticas coordinadas que reduzcan las demoras y mejoren las condiciones de operación de la red. La espera acumulada por carril se codifica mediante una transformación logarítmico-lineal acotada y discretizada, con mayor resolución para demoras bajas. Esta codificación agrupa situaciones de espera semejantes y puede facilitar el tratamiento de variaciones pequeñas, mientras limita los valores extremos. También se estudia el uso de políticas compartidas o diferenciadas ante la heterogeneidad de las intersecciones. La propuesta se evalúa como parte del sistema completo mediante simulaciones microscópicas en escenarios sintéticos y en una red vial que representa el microcentro de la Ciudad de Mendoza, comparando algoritmos independientes y cooperativos con métricas de tránsito y de aprendizaje.

\chapter*{Abstract}
\label{ch:abs}
This work addresses urban traffic congestion through a traffic-signal control system based on cooperative \acrshort{acr-marl}. It analyzes the growth of Argentina's vehicle fleet, describes the limitations of traditional signal control, and formulates the problem as a sequential decision task in which intersections act as agents. The objective is to develop and evaluate coordinated policies that reduce delays and improve network operating conditions. The system represents accumulated waiting time per lane using a bounded, discretized log-linear transformation with finer resolution for short delays. This encoding groups similar waiting situations and may facilitate the representation of small variations while limiting extreme values. The study also examines shared and differentiated policies for heterogeneous intersections. The approach is evaluated as part of the complete system through microscopic simulations of synthetic scenarios and an urban road network representing downtown Mendoza City, comparing independent and cooperative algorithms using traffic and learning metrics.
